\documentclass[a4paper,12pt,twoside]{article}
\usepackage[left=1in,right=1in,top=1in,bottom=1in]{geometry}

\usepackage{xeCJK}
\setCJKmainfont{Songti TC}
\setCJKfamilyfont{kai}{Songti TC}
\newcommand{\kai}{\CJKfamily{kai}} % 用 \kai{使用楷書}

\usepackage{lmodern}
\usepackage{titling}
\usepackage{titlesec}
\usepackage{graphicx}
\usepackage{float}
\usepackage{amsmath, amsthm, amssymb, amsfonts, mathrsfs}
\usepackage{upgreek} % 直立的希臘字母

\usepackage[shortlabels]{enumitem}
\usepackage{setspace}
\renewcommand{\baselinestretch}{1.5}
\usepackage[colorlinks]{hyperref}
\usepackage{fancyvrb}
\usepackage{fancyhdr}
\usepackage{multicol}
\usepackage{caption}
\usepackage{listings}
\usepackage{bm}
\usepackage{bbm} % for \mathbbm{1} (indicator function)
\usepackage{centernot}
\usepackage{mdframed}
\usepackage{tikz}
\usetikzlibrary{patterns, arrows.meta, positioning}
\usepackage[nameinlink, capitalise]{cleveref}
\usepackage{diagbox}
\usepackage{booktabs}
\usepackage{multirow,array}
\usepackage{tabularx}
\usepackage{mathtools}
\usepackage{indentfirst} % 首段 indent
\setlength{\parindent}{2em}
\usepackage{cancel} % 數學式子畫線槓掉
\newcommand\mcancel[2][red]{\setbox0=\hbox{$#2$}%
\rlap{\raisebox{.45\ht0}{\textcolor{#1}{\rule{\wd0}{1pt}}}}#2}
\usepackage{datetime2}

\AtBeginDocument{%
 \setlength{\abovedisplayskip}{6pt}
 \setlength{\belowdisplayskip}{6pt}
 \setlength{\abovedisplayshortskip}{4pt}
 \setlength{\belowdisplayshortskip}{4pt}
}

\definecolor{codegreen}{rgb}{0,0.6,0}
\definecolor{codegray}{rgb}{0.5,0.5,0.5}
\definecolor{codepurple}{rgb}{0.58,0,0.82}
\definecolor{backcolour}{rgb}{0.95,0.95,0.92}

\lstdefinestyle{mystyle}{
 backgroundcolor=\color{backcolour},
 commentstyle=\color{codegreen},
 keywordstyle=\color{magenta},
 numberstyle=\tiny\color{codegray},
 stringstyle=\color{codepurple},
 basicstyle=\ttfamily\footnotesize,
 breakatwhitespace=false,
 breaklines=true,
 captionpos=b,
 keepspaces=true,
 numbers=left,
 numbersep=5pt,
 showspaces=false,
 showstringspaces=false,
 showtabs=false,
 tabsize=2
}
\lstset{style=mystyle}

\hypersetup{citecolor=blue, linkcolor=blue, urlcolor=blue}

\DeclareMathOperator{\supp}{supp}
\DeclareMathOperator{\skewness}{Skewness}
\DeclareMathOperator{\DU}{D-U}
\DeclareMathOperator{\CU}{C-U}
\DeclareMathOperator{\Geo}{Geo}
\DeclareMathOperator{\Pois}{Pois}
\DeclareMathOperator{\Gam}{Gam}
\DeclareMathOperator{\LN}{LogN}
\DeclareMathOperator{\cov}{Cov}
\DeclareMathOperator*{\argmax}{arg\,max}
\DeclareMathOperator*{\argmin}{arg\,min}
\DeclareMathOperator{\re}{Re}
\DeclareMathOperator{\im}{Im}

\newcommand{\N}{\mathbb{N}}
\newcommand{\Z}{\mathbb{Z}}
\newcommand{\Q}{\mathbb{Q}}
\newcommand{\R}{\mathbb{R}}
\newcommand{\C}{\mathbb{C}}
\newcommand{\G}{\mathcal{G}}
\newcommand{\mP}{\mathbb{P}}
\newcommand{\F}{\mathcal{F}}
\newcommand{\B}{\mathcal{B}}
\newcommand{\1}{\mathbbm{1}}
\newcommand{\Poi}{\text{Poi}}
\newcommand{\Var}{\text{Var}}
\newcommand{\as}{\text{a.s.}}
\newcommand{\vect}[1]{\boldsymbol{\mathbf{#1}}}
\newcommand{\paren}[1]{\left({#1}\right)}
\newcommand{\brk}[1]{\left[{#1}\right]}
\newcommand{\curl}[1]{\left\{{#1}\right\}}
\renewcommand{\lambda}{\lambdaup}

\theoremstyle{definition} % 定理環境 style 改用正體
\newtheorem{theorem}{定理}
\newtheorem{definition}{定義}
\newtheorem*{definition*}{定義}
\newtheorem{example}{例子}
\newtheorem{exercise}{練習}
\newtheorem{corollary}{推論}
\newtheorem{property}{性質}
\newtheorem{prop}{命題}
\newtheorem{lemma}{引理}
\newtheorem{problem}{問題}
\newtheorem*{remark}{註}

\renewcommand{\figurename}{圖}
\renewcommand{\tablename}{表}

\DeclareMathOperator{\var}{Var}
\DeclareMathOperator{\corr}{corr}
\DeclareMathOperator{\rank}{rank}
\DeclareMathOperator{\MSE}{MSE}
\DeclareMathOperator{\SSE}{SSE}
\DeclareMathOperator{\SST}{SST}
\DeclareMathOperator{\SSR}{SSR}
\DeclareMathOperator{\sgn}{sign}
\DeclareMathOperator{\alt}{Alt}
\DeclareMathOperator{\adj}{adj}
\DeclareMathOperator{\curlOp}{curl}
\DeclareMathOperator{\spn}{span}
\DeclareMathOperator{\sse}{SSE}
\DeclareMathOperator{\tr}{tr}

\newcommand{\ep}{\varepsilon}
\newcommand{\p}{\partial}
\newcommand{\abs}[1]{\left|#1\right|}
\newcommand\norm[1]{\left\lVert#1\right\rVert}
\newcommand{\innerprod}[1]{\langle#1\rangle}
\newcommand{\ds}{\displaystyle}
\newcommand{\blank}{\rule[-1mm]{1.5cm}{0.1mm}}
\newcommand{\dint}{\displaystyle\int}
\renewcommand{\bar}{\overline}

\newtheoremstyle{remark}{}{}{}{}{\bfseries\sffamily}{.}{.5em}{#1}
\theoremstyle{remark}

\newtheorem*{remark*}{\textit{Remark}}
\newtheorem*{solution*}{解}

\AtEndEnvironment{solution*}{\hfill$\square$}
\newtheorem*{hint*}{Hint}


\newcommand{\redtext}[1]{{\color{red} #1}}
\newcommand{\bluetext}[1]{{\color{blue} #1}}

\fancypagestyle{firstpage}{
  \fancyhf{}
  \fancyhead[L]{\small 《高中數學實驗教材》第四冊（自然組）}
  \fancyhead[R]{\small 附錄三\quad 三角函數表}
  \fancyfoot[C]{\thepage}
}

\pagestyle{fancy}
\setlength{\headheight}{14pt}
\renewcommand{\headrulewidth}{0.5pt}
\renewcommand{\footrulewidth}{0pt}
\fancyhead{}
\fancyhead[LE]{\small 《高中數學實驗教材》第四冊（自然組）}
\fancyhead[RO]{\small 附錄三\quad 三角函數表}
\fancyfoot[C]{\thepage}

\title{\LARGE\textbf{《高中數學實驗教材》第四冊（自然組）}\\[0.5em] \Large\textbf{附錄三\quad 三角函數表}}
\author{\textbf{高中數學實驗教材編輯小組}\quad 編著}
\date{民國六十五年（1976）三版}

\begin{document}

\maketitle
\thispagestyle{firstpage}

% ==========================================
% 掃描檔勘誤、文字校訂與疑難處說明
% ==========================================
\section*{【說明】附錄三掃描檔勘誤、疑難處與文義校訂表}
\addcontentsline{toc}{section}{【說明】附錄三掃描檔勘誤、疑難處與文義校訂表}
\label{sec:errata_notes_b4_a3}

本附錄依據 1976 年三版《高中數學實驗教材 第四冊（自然組）》（高二下）附錄三「三角函數表」掃描原件（第 312～317 頁）完整轉錄為 \LaTeX{} 格式。數表內容與第三冊附錄三相同，依作業規範將掃描檔中出現之排印記號記錄彙整於下表：

\begin{table}[H]
\centering
\small
\renewcommand{\arraystretch}{1.35}
\begin{tabularx}{\textwidth}{|p{3.2cm}|p{3.8cm}|X|}
\hline
\textbf{章節位置} & \textbf{掃描檔原始現象 / 誤植} & \textbf{考證說明與校訂處理} \\ \hline

\hyperref[note:b4-a3-p314-tar]{表 2 頁 314 表頭函數記號} & 
表頭印作「$y = \text{tar }\theta$」 & 
英文正切函數記號「$\tan$」排版訛字誤植為「$\text{tar}$」。正文照原刊字樣錄入並加註說明。 \\ \hline

\end{tabularx}
\end{table}

\clearpage
\tableofcontents
\newpage

\section*{附錄三\quad 三角函數表}
\addcontentsline{toc}{section}{附錄三\quad 三角函數表}
\label{sec:b4_a3}

\subsection*{表 1.\quad 三角函數表（正弦餘弦）}
\addcontentsline{toc}{subsection}{表 1.\quad 三角函數表（正弦餘弦）}
\begin{center} $y = \sin\theta$ \end{center}
{\centering\fontsize{7.5pt}{8.8pt}\selectfont\setlength{\tabcolsep}{3.5pt}\renewcommand{\arraystretch}{0.96}
\begin{tabular}{|c|cccccccccc|}
\hline
$\theta\text{ (degrees)}$ & 0 & 0.1 & 0.2 & 0.3 & 0.4 & 0.5 & 0.6 & 0.7 & 0.8 & 0.9 \\ \hline
0 & 0.0000 & 0.0017 & 0.0035 & 0.0052 & 0.0070 & 0.0087 & 0.0105 & 0.0122 & 0.0140 & 0.0157 \\
1 & 0.0175 & 0.0192 & 0.0209 & 0.0227 & 0.0244 & 0.0262 & 0.0279 & 0.0297 & 0.0314 & 0.0332 \\
2 & 0.0349 & 0.0366 & 0.0384 & 0.0401 & 0.0419 & 0.0436 & 0.0454 & 0.0471 & 0.0488 & 0.0506 \\
3 & 0.0523 & 0.0541 & 0.0558 & 0.0576 & 0.0593 & 0.0610 & 0.0628 & 0.0645 & 0.0663 & 0.0680 \\
4 & 0.0698 & 0.0715 & 0.0732 & 0.0750 & 0.0767 & 0.0785 & 0.0802 & 0.0819 & 0.0837 & 0.0854 \\
\hline
5 & 0.0812 & 0.0889 & 0.0906 & 0.0924 & 0.0941 & 0.0958 & 0.0976 & 0.0993 & 0.1011 & 0.1028 \\
6 & 0.1045 & 0.1063 & 0.1080 & 0.1097 & 0.1115 & 0.1132 & 0.1149 & 0.1167 & 0.1184 & 0.1201 \\
7 & 0.1219 & 0.1236 & 0.1253 & 0.1271 & 0.1288 & 0.1305 & 0.1323 & 0.1340 & 0.1357 & 0.1374 \\
8 & 0.1392 & 0.1409 & 0.1426 & 0.1444 & 0.1461 & 0.1478 & 0.1495 & 0.1513 & 0.1530 & 0.1547 \\
9 & 0.1564 & 0.1582 & 0.1599 & 0.1616 & 0.1633 & 0.1650 & 0.1668 & 0.1685 & 0.1702 & 0.1719 \\
\hline
10 & 0.1736 & 0.1754 & 0.1771 & 0.1788 & 0.1805 & 0.1822 & 0.1839 & 0.1857 & 0.1874 & 0.1891 \\
11 & 0.1908 & 0.1925 & 0.1942 & 0.1959 & 0.1977 & 0.1994 & 0.2011 & 0.2028 & 0.2045 & 0.2062 \\
12 & 0.2079 & 0.2096 & 0.2113 & 0.2130 & 0.2147 & 0.2164 & 0.2181 & 0.2198 & 0.2215 & 0.2232 \\
13 & 0.2249 & 0.2266 & 0.2283 & 0.2300 & 0.2317 & 0.2334 & 0.2351 & 0.2368 & 0.2385 & 0.2402 \\
14 & 0.2419 & 0.2436 & 0.2453 & 0.2470 & 0.2487 & 0.2504 & 0.2521 & 0.2538 & 0.2554 & 0.2571 \\
\hline
15 & 0.2588 & 0.2605 & 0.2622 & 0.2639 & 0.2656 & 0.2672 & 0.2689 & 0.2706 & 0.2723 & 0.2740 \\
16 & 0.2756 & 0.2773 & 0.2790 & 0.2807 & 0.2823 & 0.2840 & 0.2857 & 0.2874 & 0.2890 & 0.2907 \\
17 & 0.2924 & 0.2940 & 0.2957 & 0.2974 & 0.2990 & 0.3007 & 0.3024 & 0.3040 & 0.3057 & 0.3074 \\
18 & 0.3090 & 0.3107 & 0.3123 & 0.3140 & 0.3156 & 0.3173 & 0.3190 & 0.3206 & 0.3223 & 0.3239 \\
19 & 0.3256 & 0.3272 & 0.3289 & 0.3305 & 0.3322 & 0.3338 & 0.3354 & 0.3371 & 0.3387 & 0.3404 \\
\hline
20 & 0.3420 & 0.3437 & 0.3453 & 0.3469 & 0.3486 & 0.3502 & 0.3518 & 0.3535 & 0.3551 & 0.3567 \\
21 & 0.3584 & 0.3600 & 0.3616 & 0.3632 & 0.3649 & 0.3665 & 0.3681 & 0.3697 & 0.3714 & 0.3730 \\
22 & 0.3746 & 0.3762 & 0.3778 & 0.3795 & 0.3811 & 0.3827 & 0.3843 & 0.3859 & 0.3875 & 0.3891 \\
23 & 0.3907 & 0.3923 & 0.3939 & 0.3955 & 0.3971 & 0.3987 & 0.4003 & 0.4019 & 0.4035 & 0.4051 \\
24 & 0.4067 & 0.4083 & 0.4099 & 0.4115 & 0.4131 & 0.4147 & 0.4163 & 0.4179 & 0.4194 & 0.4210 \\
\hline
25 & 0.4226 & 0.4242 & 0.4258 & 0.4274 & 0.4289 & 0.4305 & 0.4321 & 0.4337 & 0.4352 & 0.4368 \\
26 & 0.4384 & 0.4399 & 0.4415 & 0.4431 & 0.4446 & 0.4462 & 0.4478 & 0.4493 & 0.4509 & 0.4524 \\
27 & 0.4540 & 0.4555 & 0.4571 & 0.4586 & 0.4602 & 0.4617 & 0.4633 & 0.4648 & 0.4664 & 0.4679 \\
28 & 0.4695 & 0.4710 & 0.4725 & 0.4741 & 0.4756 & 0.4772 & 0.4787 & 0.4802 & 0.4817 & 0.4833 \\
29 & 0.4848 & 0.4863 & 0.4879 & 0.4894 & 0.4909 & 0.4924 & 0.4939 & 0.4955 & 0.4970 & 0.4985 \\
\hline
30 & 0.5000 & 0.5015 & 0.5030 & 0.5045 & 0.5060 & 0.5075 & 0.5090 & 0.5105 & 0.5120 & 0.5135 \\
31 & 0.5150 & 0.5165 & 0.5180 & 0.5195 & 0.5210 & 0.5225 & 0.5240 & 0.5255 & 0.5269 & 0.5284 \\
32 & 0.5299 & 0.5314 & 0.5329 & 0.5344 & 0.5358 & 0.5373 & 0.5388 & 0.5402 & 0.5417 & 0.5432 \\
33 & 0.5446 & 0.5461 & 0.5476 & 0.5490 & 0.5505 & 0.5519 & 0.5534 & 0.5548 & 0.5563 & 0.5577 \\
34 & 0.5592 & 0.5606 & 0.5621 & 0.5635 & 0.5650 & 0.5664 & 0.5678 & 0.5693 & 0.5707 & 0.5721 \\
\hline
35 & 0.5736 & 0.5750 & 0.5764 & 0.5778 & 0.5793 & 0.5807 & 0.5821 & 0.5835 & 0.5849 & 0.5864 \\
36 & 0.5878 & 0.5892 & 0.5906 & 0.5920 & 0.5934 & 0.5948 & 0.5962 & 0.5976 & 0.5990 & 0.6004 \\
37 & 0.6018 & 0.6032 & 0.6046 & 0.6060 & 0.6074 & 0.6088 & 0.6101 & 0.6115 & 0.6129 & 0.6143 \\
38 & 0.6157 & 0.6170 & 0.6184 & 0.6198 & 0.6211 & 0.6225 & 0.6239 & 0.6252 & 0.6266 & 0.6280 \\
39 & 0.6293 & 0.6307 & 0.6320 & 0.6334 & 0.6347 & 0.6361 & 0.6374 & 0.6388 & 0.6401 & 0.6414 \\
\hline
40 & 0.6428 & 0.6441 & 0.6454 & 0.6468 & 0.6481 & 0.6494 & 0.6508 & 0.6521 & 0.6534 & 0.6547 \\
41 & 0.6560 & 0.6574 & 0.6587 & 0.6600 & 0.6613 & 0.6626 & 0.6639 & 0.6652 & 0.6665 & 0.6678 \\
42 & 0.6691 & 0.6704 & 0.6717 & 0.6730 & 0.6743 & 0.6756 & 0.6769 & 0.6782 & 0.6794 & 0.6807 \\
43 & 0.6820 & 0.6833 & 0.6845 & 0.6858 & 0.6871 & 0.6883 & 0.6896 & 0.6909 & 0.6921 & 0.6934 \\
44 & 0.6946 & 0.6959 & 0.6972 & 0.6984 & 0.6997 & 0.7009 & 0.7021 & 0.7034 & 0.7046 & 0.7059 \\
\hline
$\theta\text{ (degrees)}$ & 0 & 0.1 & 0.2 & 0.3 & 0.4 & 0.5 & 0.6 & 0.7 & 0.8 & 0.9 \\ \hline
\end{tabular}\par}
\vspace{0.5em}\begin{center}\small 餘弦 $\cos\theta = \sin(90^\circ - \theta)$\end{center}
\clearpage

\subsection*{表 1.\quad （續）三角函數表（正弦餘弦）}
\begin{center} $y = \sin\theta$ \end{center}
{\centering\fontsize{7.5pt}{8.8pt}\selectfont\setlength{\tabcolsep}{3.5pt}\renewcommand{\arraystretch}{0.96}
\begin{tabular}{|c|cccccccccc|}
\hline
$\theta\text{ (degrees)}$ & 0 & 0.1 & 0.2 & 0.3 & 0.4 & 0.5 & 0.6 & 0.7 & 0.8 & 0.9 \\ \hline
45 & 0.7071 & 0.7083 & 0.7096 & 0.7108 & 0.7120 & 0.7133 & 0.7145 & 0.7157 & 0.7169 & 0.7181 \\
46 & 0.7083 & 0.7096 & 0.7218 & 0.7120 & 0.7132 & 0.7145 & 0.7266 & 0.7169 & 0.7181 & 0.7302 \\
47 & 0.7205 & 0.7218 & 0.7230 & 0.7242 & 0.7254 & 0.7266 & 0.7385 & 0.7290 & 0.7302 & 0.7420 \\
48 & 0.7325 & 0.7337 & 0.7349 & 0.7466 & 0.7373 & 0.7384 & 0.7396 & 0.7408 & 0.7420 & 0.7536 \\
49 & 0.7547 & 0.7455 & 0.7466 & 0.7581 & 0.7489 & 0.7604 & 0.7513 & 0.7524 & 0.7536 & 0.7649 \\
\hline
50 & 0.7558 & 0.7570 & 0.7581 & 0.7593 & 0.7604 & 0.7615 & 0.7727 & 0.7638 & 0.7649 & 0.7760 \\
51 & 0.7672 & 0.7683 & 0.7694 & 0.7705 & 0.7716 & 0.7727 & 0.7738 & 0.7848 & 0.7760 & 0.7869 \\
52 & 0.7880 & 0.7793 & 0.7902 & 0.7815 & 0.7826 & 0.7837 & 0.7944 & 0.7858 & 0.7869 & 0.7976 \\
53 & 0.7891 & 0.7901 & 0.7912 & 0.7923 & 0.8028 & 0.7944 & 0.7955 & 0.7965 & 0.7976 & 0.8080 \\
54 & 0.7997 & 0.8007 & 0.8018 & 0.8028 & 0.8038 & 0.8049 & 0.8059 & 0.8161 & 0.8080 & 0.8181 \\
\hline
55 & 0.8100 & 0.8111 & 0.8121 & 0.8131 & 0.8231 & 0.8241 & 0.8251 & 0.8261 & 0.8181 & 0.8281 \\
56 & 0.8201 & 0.8211 & 0.8221 & 0.8231 & 0.8241 & 0.8339 & 0.8348 & 0.8271 & 0.8281 & 0.8377 \\
57 & 0.8300 & 0.8310 & 0.8319 & 0.8329 & 0.8339 & 0.8348 & 0.8358 & 0.8368 & 0.8377 & 0.8471 \\
58 & 0.8396 & 0.8406 & 0.8415 & 0.8424 & 0.8434 & 0.8443 & 0.8453 & 0.8462 & 0.8471 & 0.8563 \\
59 & 0.8490 & 0.8581 & 0.8508 & 0.8517 & 0.8526 & 0.8535 & 0.8544 & 0.8554 & 0.8563 & 0.8652 \\
\hline
60 & 0.8581 & 0.8590 & 0.8598 & 0.8607 & 0.8616 & 0.8625 & 0.8712 & 0.8643 & 0.8729 & 0.8738 \\
61 & 0.8669 & 0.8678 & 0.8686 & 0.8695 & 0.8780 & 0.8788 & 0.8721 & 0.8729 & 0.8738 & 0.8821 \\
62 & 0.8755 & 0.8763 & 0.8771 & 0.8780 & 0.8788 & 0.8796 & 0.8878 & 0.8886 & 0.8821 & 0.8902 \\
63 & 0.8838 & 0.8846 & 0.8854 & 0.8934 & 0.8870 & 0.8878 & 0.8957 & 0.8894 & 0.8902 & 0.8980 \\
64 & 0.8918 & 0.8926 & 0.8934 & 0.8941 & 0.8949 & 0.8957 & 0.8965 & 0.8973 & 0.8980 & 0.9056 \\
\hline
65 & 0.8995 & 0.9003 & 0.9011 & 0.9018 & 0.9026 & 0.9033 & 0.9041 & 0.9114 & 0.9056 & 0.9128 \\
66 & 0.9070 & 0.9078 & 0.9085 & 0.9092 & 0.9100 & 0.9107 & 0.9178 & 0.9121 & 0.9128 & 0.9198 \\
67 & 0.9142 & 0.9150 & 0.9157 & 0.9164 & 0.9171 & 0.9239 & 0.9245 & 0.9191 & 0.9198 & 0.9265 \\
68 & 0.9212 & 0.9219 & 0.9225 & 0.9232 & 0.9239 & 0.9245 & 0.9252 & 0.9259 & 0.9265 & 0.9330 \\
69 & 0.9278 & 0.9285 & 0.9291 & 0.9298 & 0.9304 & 0.9367 & 0.9317 & 0.9379 & 0.9329 & 0.9391 \\
\hline
70 & 0.9342 & 0.9348 & 0.9354 & 0.9361 & 0.9367 & 0.9373 & 0.9379 & 0.9438 & 0.9444 & 0.9449 \\
71 & 0.9403 & 0.9409 & 0.9415 & 0.9421 & 0.9478 & 0.9432 & 0.9489 & 0.9444 & 0.9449 & 0.9505 \\
72 & 0.9461 & 0.9466 & 0.9472 & 0.9478 & 0.9483 & 0.9489 & 0.9542 & 0.9500 & 0.9505 & 0.9558 \\
73 & 0.9516 & 0.9521 & 0.9527 & 0.9532 & 0.9537 & 0.9542 & 0.9548 & 0.9553 & 0.9558 & 0.9608 \\
74 & 0.9568 & 0.9573 & 0.9578 & 0.9583 & 0.9588 & 0.9593 & 0.9598 & 0.9603 & 0.9650 & 0.9655 \\
\hline
75 & 0.9617 & 0.9622 & 0.9627 & 0.9632 & 0.9636 & 0.9681 & 0.9646 & 0.9650 & 0.9694 & 0.9699 \\
76 & 0.9664 & 0.9668 & 0.9673 & 0.9677 & 0.9651 & 0.9686 & 0.9728 & 0.9694 & 0.9699 & 0.9740 \\
77 & 0.9707 & 0.9711 & 0.9715 & 0.9755 & 0.9724 & 0.9728 & 0.9732 & 0.9770 & 0.9774 & 0.9778 \\
78 & 0.9748 & 0.9751 & 0.9755 & 0.9759 & 0.9763 & 0.9799 & 0.9770 & 0.9806 & 0.9810 & 0.9813 \\
79 & 0.9785 & 0.9789 & 0.9792 & 0.9796 & 0.9799 & 0.9833 & 0.9806 & 0.9839 & 0.9842 & 0.9845 \\
\hline
80 & 0.9848 & 0.9823 & 0.9826 & 0.9857 & 0.9833 & 0.9836 & 0.9539 & 0.9869 & 0.9843 & 0.9874 \\
81 & 0.9877 & 0.9854 & 0.9857 & 0.9860 & 0.9863 & 0.9890 & 0.9893 & 0.9871 & 0.9874 & 0.9900 \\
82 & 0.9880 & 0.9882 & 0.9885 & 0.9888 & 0.9890 & 0.9893 & 0.9917 & 0.9919 & 0.9900 & 0.9923 \\
83 & 0.9905 & 0.9907 & 0.9910 & 0.9932 & 0.9914 & 0.9936 & 0.9938 & 0.9921 & 0.9942 & 0.9943 \\
84 & 0.9928 & 0.9930 & 0.9932 & 0.9934 & 0.9936 & 0.9938 & 0.9956 & 0.9957 & 0.9943 & 0.9960 \\
\hline
85 & 0.9947 & 0.9949 & 0.9951 & 0.9952 & 0.9954 & 0.9956 & 0.9957 & 0.9959 & 0.9960 & 0.9974 \\
86 & 0.9963 & 0.9965 & 0.9966 & 0.9968 & 0.9980 & 0.9971 & 0.9982 & 0.9983 & 0.9974 & 0.9985 \\
87 & 0.9977 & 0.9978 & 0.9979 & 0.9980 & 0.9981 & 0.9982 & 0.9983 & 0.9984 & 0.9993 & 0.9993 \\
88 & 0.9987 & 0.9988 & 0.9989 & 0.9990 & 0.9990 & 0.9991 & 0.9992 & 0.9993 & 0.9993 & 0.9998 \\
89 & 0.9994 & 0.9995 & 0.9996 & 0.9996 & 0.9997 & 0.9997 & 0.9997 & 0.9998 & 0.9998 & 1.0000 \\
\hline
$\theta\text{ (degrees)}$ & 0 & 0.1 & 0.2 & 0.3 & 0.4 & 0.5 & 0.6 & 0.7 & 0.8 & 0.9 \\ \hline
\end{tabular}\par}
\vspace{0.5em}\begin{center}\small 餘弦 $\cos\theta = \sin(90^\circ - \theta)$\end{center}
\clearpage

\subsection*{表 2.\quad 三角函數表（正切餘切）}
\addcontentsline{toc}{subsection}{表 2.\quad 三角函數表（正切餘切）}
\begin{center} $y = \text{tar }\theta$\phantomsection\label{note:b4-a3-p314-tar}\footnote{【原書註記/校訂】：原書頁 314 此處表頭印作「$y = \text{tar }\theta$」，實為「$y = \tan\theta$」之排印訛字．正文照原刊字樣錄入．} \end{center}
{\centering\fontsize{7.5pt}{8.8pt}\selectfont\setlength{\tabcolsep}{3.5pt}\renewcommand{\arraystretch}{0.96}
\begin{tabular}{|c|cccccccccc|}
\hline
$\theta\text{ (degrees)}$ & 0 & 0.1 & 0.2 & 0.3 & 0.4 & 0.5 & 0.6 & 0.7 & 0.8 & 0.9 \\ \hline
0 & 0.0000 & 0.0017 & 0.0035 & 0.0052 & 0.0070 & 0.0087 & 0.0105 & 0.0122 & 0.0140 & 0.0157 \\
1 & 0.0000 & 0.0017 & 0.0035 & 0.0052 & 0.0070 & 0.0087 & 0.0105 & 0.0122 & 0.0140 & 0.0157 \\
2 & 0.0175 & 0.0192 & 0.0209 & 0.0227 & 0.0244 & 0.0262 & 0.0279 & 0.0297 & 0.0314 & 0.0332 \\
3 & 0.0349 & 0.0542 & 0.0384 & 0.0402 & 0.0419 & 0.0612 & 0.0454 & 0.0472 & 0.0489 & 0.0507 \\
4 & 0.0524 & 0.0542 & 0.0559 & 0.0577 & 0.0594 & 0.0612 & 0.0629 & 0.0647 & 0.0664 & 0.0682 \\
\hline
5 & 0.0875 & 0.0717 & 0.0734 & 0.0752 & 0.0769 & 0.0787 & 0.0805 & 0.0822 & 0.0840 & 0.0857 \\
6 & 0.0875 & 0.0892 & 0.0910 & 0.0928 & 0.0945 & 0.0963 & 0.0980 & 0.0998 & 0.1016 & 0.1033 \\
7 & 0.1051 & 0.1069 & 0.1086 & 0.1104 & 0.1122 & 0.1139 & 0.1157 & 0.1175 & 0.1370 & 0.1210 \\
8 & 0.1228 & 0.1246 & 0.1263 & 0.1281 & 0.1299 & 0.1317 & 0.1334 & 0.1530 & 0.1370 & 0.1388 \\
9 & 0.1405 & 0.1423 & 0.1441 & 0.1459 & 0.1477 & 0.1494 & 0.1512 & 0.1530 & 0.1548 & 0.1566 \\
\hline
10 & 0.1584 & 0.1602 & 0.1620 & 0.1638 & 0.1655 & 0.1673 & 0.1691 & 0.1709 & 0.1727 & 0.1745 \\
11 & 0.1763 & 0.1962 & 0.1799 & 0.1817 & 0.1835 & 0.1853 & 0.1871 & 0.1889 & 0.1908 & 0.1926 \\
12 & 0.1944 & 0.1962 & 0.1980 & 0.1998 & 0.2016 & 0.2034 & 0.2053 & 0.2071 & 0.2089 & 0.2107 \\
13 & 0.2126 & 0.2327 & 0.2162 & 0.2180 & 0.2199 & 0.2217 & 0.2235 & 0.2254 & 0.2272 & 0.2290 \\
14 & 0.2309 & 0.2327 & 0.2345 & 0.2364 & 0.2382 & 0.2401 & 0.2419 & 0.2438 & 0.2456 & 0.2475 \\
\hline
15 & 0.2493 & 0.2512 & 0.2530 & 0.2549 & 0.2568 & 0.2586 & 0.2605 & 0.2623 & 0.2642 & 0.2661 \\
16 & 0.2679 & 0.2698 & 0.2717 & 0.2736 & 0.2754 & 0.2773 & 0.2792 & 0.2811 & 0.2830 & 0.2849 \\
17 & 0.2867 & 0.2886 & 0.2905 & 0.2924 & 0.2945 & 0.2962 & 0.2981 & 0.3000 & 0.3019 & 0.3038 \\
18 & 0.3057 & 0.3076 & 0.3095 & 0.3115 & 0.3134 & 0.3153 & 0.3172 & 0.3191 & 0.3211 & 0.3230 \\
19 & 0.3249 & 0.3268 & 0.3288 & 0.3307 & 0.3326 & 0.3346 & 0.3365 & 0.3385 & 0.3404 & 0.3424 \\
\hline
20 & 0.3443 & 0.3463 & 0.3482 & 0.3502 & 0.3521 & 0.3541 & 0.3561 & 0.3580 & 0.3600 & 0.3620 \\
21 & 0.3640 & 0.3659 & 0.3679 & 0.3699 & 0.3719 & 0.3739 & 0.3759 & 0.3779 & 0.3799 & 0.3819 \\
22 & 0.3839 & 0.3859 & 0.3879 & 0.3899 & 0.3919 & 0.3939 & 0.3959 & 0.3979 & 0.4000 & 0.4020 \\
23 & 0.4040 & 0.4061 & 0.4081 & 0.4101 & 0.4122 & 0.4142 & 0.4163 & 0.4183 & 0.4204 & 0.4224 \\
24 & 0.4245 & 0.4265 & 0.4286 & 0.4307 & 0.4327 & 0.4348 & 0.4369 & 0.4390 & 0.4410 & 0.4431 \\
\hline
25 & 0.4452 & 0.4473 & 0.4494 & 0.4515 & 0.4536 & 0.4557 & 0.4578 & 0.4599 & 0.4621 & 0.4642 \\
26 & 0.4663 & 0.4684 & 0.4706 & 0.4727 & 0.4748 & 0.4770 & 0.4791 & 0.4813 & 0.4834 & 0.4856 \\
27 & 0.4877 & 0.4899 & 0.4921 & 0.4942 & 0.4964 & 0.4986 & 0.5008 & 0.5029 & 0.5051 & 0.5073 \\
28 & 0.5095 & 0.5117 & 0.5139 & 0.5161 & 0.5183 & 0.5206 & 0.5228 & 0.5250 & 0.5272 & 0.5295 \\
29 & 0.5317 & 0.5339 & 0.5362 & 0.5384 & 0.5407 & 0.5429 & 0.5452 & 0.5475 & 0.5497 & 0.5520 \\
\hline
30 & 0.5774 & 0.5566 & 0.5589 & 0.5612 & 0.5635 & 0.5658 & 0.5681 & 0.5704 & 0.5727 & 0.5750 \\
31 & 0.5773 & 0.5797 & 0.5820 & 0.5843 & 0.5867 & 0.5890 & 0.5914 & 0.5937 & 0.5961 & 0.5985 \\
32 & 0.6008 & 0.6032 & 0.6056 & 0.6080 & 0.6104 & 0.6128 & 0.6152 & 0.6176 & 0.6200 & 0.6224 \\
33 & 0.6249 & 0.6273 & 0.6297 & 0.6322 & 0.6346 & 0.6371 & 0.6395 & 0.6420 & 0.6444 & 0.6469 \\
34 & 0.6494 & 0.6519 & 0.6544 & 0.6569 & 0.6594 & 0.6619 & 0.6644 & 0.6669 & 0.6694 & 0.6720 \\
\hline
35 & 0.6745 & 0.6770 & 0.6796 & 0.6821 & 0.6847 & 0.6873 & 0.6898 & 0.6924 & 0.6950 & 0.6976 \\
36 & 0.7002 & 0.7028 & 0.7054 & 0.7080 & 0.7106 & 0.7133 & 0.7159 & 0.7186 & 0.7212 & 0.7239 \\
37 & 0.7265 & 0.7292 & 0.7319 & 0.7346 & 0.7372 & 0.7399 & 0.7701 & 0.7454 & 0.7481 & 0.7508 \\
38 & 0.7535 & 0.7563 & 0.7590 & 0.7618 & 0.7645 & 0.7673 & 0.7701 & 0.7729 & 0.7757 & 0.7785 \\
39 & 0.7813 & 0.7841 & 0.7869 & 0.7897 & 0.7926 & 0.7954 & 0.7983 & 0.8011 & 0.8040 & 0.8069 \\
\hline
40 & 0.8098 & 0.8127 & 0.8156 & 0.8185 & 0.8214 & 0.8243 & 0.8273 & 0.8302 & 0.8331 & 0.8361 \\
41 & 0.8391 & 0.8421 & 0.8450 & 0.8480 & 0.8510 & 0.8541 & 0.8571 & 0.8601 & 0.8632 & 0.8662 \\
42 & 0.8693 & 0.8723 & 0.8754 & 0.8785 & 0.8816 & 0.8847 & 0.8878 & 0.8909 & 0.8941 & 0.8972 \\
43 & 0.9004 & 0.9035 & 0.9067 & 0.9099 & 0.9131 & 0.9163 & 0.9195 & 0.9228 & 0.9260 & 0.9292 \\
44 & 0.9325 & 0.9358 & 0.9390 & 0.9423 & 0.9456 & 0.9489 & 0.9523 & 0.9556 & 0.9589 & 0.9623 \\
\hline
$\theta\text{ (degrees)}$ & 0 & 0.1 & 0.2 & 0.3 & 0.4 & 0.5 & 0.6 & 0.7 & 0.8 & 0.9 \\ \hline
\end{tabular}\par}
\vspace{0.5em}\begin{center}\small 餘切 $\cot\theta = \tan(90^\circ - \theta)$\end{center}
\clearpage

\subsection*{表 2.\quad （續）三角函數表（正切餘切）}
\begin{center} $y = \tan\theta$ \end{center}
{\centering\fontsize{7.5pt}{8.8pt}\selectfont\setlength{\tabcolsep}{3.5pt}\renewcommand{\arraystretch}{0.96}
\begin{tabular}{|c|cccccccccc|}
\hline
$\theta\text{ (degrees)}$ & 0 & 0.1 & 0.2 & 0.3 & 0.4 & 0.5 & 0.6 & 0.7 & 0.8 & 0.9 \\ \hline
45 & 1.0000 & 1.0035 & 1.0070 & 1.0105 & 1.0141 & 1.0176 & 1.0212 & 1.0247 & 1.0283 & 1.0319 \\
46 & 1.0355 & 1.0070 & 1.0105 & 1.0140 & 1.0176 & 1.0211 & 1.0247 & 1.0283 & 1.0319 & 1.0686 \\
47 & 1.0391 & 1.0428 & 1.0464 & 1.0501 & 1.0538 & 1.0913 & 1.0951 & 1.0649 & 1.0686 & 1.1067 \\
48 & 1.0761 & 1.0799 & 1.0837 & 1.0875 & 1.0913 & 1.1303 & 1.0990 & 1.1028 & 1.1067 & 1.1463 \\
49 & 1.1145 & 1.1184 & 1.1223 & 1.1263 & 1.1303 & 1.1708 & 1.1382 & 1.1423 & 1.1463 & 1.1875 \\
\hline
50 & 1.1918 & 1.1585 & 1.1626 & 1.1667 & 1.1708 & 1.1750 & 1.1791 & 1.1833 & 1.1875 & 1.2305 \\
51 & 1.2349 & 1.2002 & 1.2045 & 1.2088 & 1.2131 & 1.2572 & 1.2217 & 1.2261 & 1.2305 & 1.2753 \\
52 & 1.2393 & 1.2437 & 1.2892 & 1.2526 & 1.2985 & 1.2616 & 1.3079 & 1.2707 & 1.2753 & 1.3222 \\
53 & 1.3270 & 1.2892 & 1.2938 & 1.2985 & 1.3032 & 1.3514 & 1.3126 & 1.3174 & 1.3222 & 1.3713 \\
54 & 1.3318 & 1.3367 & 1.3865 & 1.3465 & 1.3514 & 1.3563 & 1.3613 & 1.3663 & 1.3713 & 1.4229 \\
\hline
55 & 1.3814 & 1.3865 & 1.3916 & 1.3967 & 1.4019 & 1.4071 & 1.4123 & 1.4175 & 1.4715 & 1.4770 \\
56 & 1.4334 & 1.4388 & 1.4441 & 1.4495 & 1.5051 & 1.5108 & 1.5166 & 1.5224 & 1.5282 & 1.5340 \\
57 & 1.5399 & 1.5458 & 1.5517 & 1.5577 & 1.5637 & 1.5697 & 1.5757 & 1.5818 & 1.5880 & 1.5941 \\
58 & 1.6003 & 1.6066 & 1.6128 & 1.6191 & 1.6255 & 1.6319 & 1.6383 & 1.6447 & 1.6512 & 1.6577 \\
59 & 1.6643 & 1.6709 & 1.6775 & 1.6842 & 1.6909 & 1.6977 & 1.7045 & 1.7113 & 1.7182 & 1.7251 \\
\hline
60 & 1.7321 & 1.7391 & 1.7461 & 1.7532 & 1.7603 & 1.7675 & 1.7747 & 1.7820 & 1.7893 & 1.7966 \\
61 & 1.8040 & 1.8115 & 1.8190 & 1.8265 & 1.8341 & 1.8418 & 1.8495 & 1.8572 & 1.8650 & 1.8728 \\
62 & 1.8807 & 1.8887 & 1.8967 & 1.9047 & 1.9128 & 1.9210 & 1.9292 & 1.9375 & 1.9458 & 1.9542 \\
63 & 1.9626 & 1.9711 & 1.9797 & 1.9883 & 1.9970 & 2.0057 & 2.0145 & 2.0233 & 2.0323 & 2.0413 \\
64 & 2.0503 & 2.0594 & 2.0686 & 2.0778 & 2.0872 & 2.0965 & 2.1060 & 2.1155 & 2.1251 & 2.1348 \\
\hline
65 & 2.1445 & 2.1543 & 2.1642 & 2.1742 & 2.1842 & 2.1943 & 2.2045 & 2.2148 & 2.2251 & 2.2355 \\
66 & 2.2460 & 2.2566 & 2.2673 & 2.2781 & 2.2889 & 2.2998 & 2.3109 & 2.3220 & 2.3332 & 2.3445 \\
67 & 2.3559 & 2.3673 & 2.3789 & 2.3906 & 2.4023 & 2.4142 & 2.4262 & 2.4383 & 2.4504 & 2.4627 \\
68 & 2.4751 & 2.4876 & 2.5002 & 2.5129 & 2.5257 & 2.5386 & 2.5517 & 2.5649 & 2.5782 & 2.5916 \\
69 & 2.6051 & 2.6187 & 2.6325 & 2.6464 & 2.6605 & 2.6746 & 2.6889 & 2.7034 & 2.7179 & 2.7326 \\
\hline
70 & 2.7475 & 2.7625 & 2.7776 & 2.7929 & 2.8083 & 2.8239 & 2.8397 & 2.8556 & 2.8716 & 2.8878 \\
71 & 2.9042 & 2.9208 & 2.9375 & 2.9544 & 2.9714 & 2.9887 & 3.0061 & 3.0237 & 3.0415 & 3.0595 \\
72 & 3.0777 & 3.0961 & 3.1146 & 3.1334 & 3.1524 & 3.1716 & 3.1910 & 3.2106 & 3.2305 & 3.2506 \\
73 & 3.2709 & 3.2914 & 3.3122 & 3.3332 & 3.3544 & 3.3759 & 3.3977 & 3.4197 & 3.4420 & 3.4646 \\
74 & 3.4874 & 3.5105 & 3.5339 & 3.5576 & 3.5816 & 3.6059 & 3.6305 & 3.6554 & 3.6806 & 3.7062 \\
\hline
75 & 3.7321 & 3.7583 & 3.7848 & 3.8118 & 3.8391 & 3.8667 & 3.8947 & 3.9232 & 3.9520 & 3.9812 \\
76 & 4.0108 & 4.0408 & 4.0713 & 4.1022 & 4.1335 & 4.1653 & 4.1976 & 4.2303 & 4.2635 & 4.2972 \\
77 & 4.3315 & 4.3662 & 4.4015 & 4.4373 & 4.4737 & 4.5107 & 4.5483 & 4.5864 & 4.6252 & 4.6646 \\
78 & 4.7046 & 4.7453 & 4.7867 & 4.8288 & 4.8716 & 4.9152 & 4.9594 & 5.0045 & 5.0504 & 5.0970 \\
79 & 5.1446 & 5.1929 & 5.2422 & 5.2924 & 5.3435 & 5.3955 & 5.4486 & 5.5026 & 5.5578 & 5.6140 \\
\hline
80 & 5.6713 & 5.7297 & 5.7894 & 5.8502 & 5.9124 & 5.9758 & 6.0405 & 6.1066 & 6.1742 & 6.2432 \\
81 & 6.3138 & 6.3859 & 6.4596 & 6.5350 & 6.6122 & 6.6912 & 6.7720 & 6.8548 & 6.9395 & 7.0264 \\
82 & 7.1154 & 7.2066 & 7.3002 & 7.3962 & 7.4947 & 7.5958 & 7.6996 & 7.8062 & 7.9158 & 8.0285 \\
83 & 8.1443 & 8.2636 & 8.3863 & 8.5126 & 8.6427 & 8.7769 & 8.9152 & 9.0579 & 9.2052 & 9.3572 \\
84 & 9.5144 & 9.6768 & 9.8448 & 10.0187 & 10.1988 & 10.3854 & 10.5789 & 10.7797 & 10.9882 & 11.2048 \\
\hline
85 & 11.4301 & 11.6645 & 11.9087 & 12.1632 & 12.4288 & 12.7062 & 12.9962 & 13.2996 & 13.6174 & 13.9507 \\
86 & 14.3007 & 14.6685 & 15.0557 & 15.4638 & 15.8945 & 16.3499 & 16.8319 & 17.3432 & 17.8863 & 18.4645 \\
87 & 19.0811 & 19.7403 & 20.4465 & 21.2049 & 22.0217 & 22.9038 & 23.8593 & 24.8978 & 26.0307 & 27.2715 \\
88 & 28.6363 & 30.1446 & 31.8205 & 33.6935 & 35.8006 & 38.1885 & 40.9174 & 44.0661 & 47.7395 & 52.0807 \\
\hline
$\theta\text{ (degrees)}$ & 0 & 0.1 & 0.2 & 0.3 & 0.4 & 0.5 & 0.6 & 0.7 & 0.8 & 0.9 \\ \hline
\end{tabular}\par}
\vspace{0.5em}\begin{center}\small 餘切 $\cot\theta = \tan(90^\circ - \theta)$\end{center}
\clearpage

\subsection*{表 3.\quad $360^\circ$ 與弧度換算表}
\addcontentsline{toc}{subsection}{表 3.\quad $360^\circ$ 與弧度換算表}
{\centering\fontsize{7.5pt}{8.8pt}\selectfont\setlength{\tabcolsep}{3.5pt}\renewcommand{\arraystretch}{0.96}
\begin{tabular}{|c|cccccccccc|}
\hline
$\theta\text{ (degrees)}$ & 0 & 0.1 & 0.2 & 0.3 & 0.4 & 0.5 & 0.6 & 0.7 & 0.8 & 0.9 \\ \hline
0 & 0.0000 & 0.0017 & 0.0035 & 0.0052 & 0.0070 & 0.0087 & 0.0105 & 0.0122 & 0.0140 & 0.0157 \\
1 & 0.0000 & 0.0017 & 0.0035 & 0.0052 & 0.0244 & 0.0087 & 0.0105 & 0.0122 & 0.0140 & 0.0157 \\
2 & 0.0175 & 0.0192 & 0.0209 & 0.0227 & 0.0244 & 0.0262 & 0.0279 & 0.0297 & 0.0314 & 0.0332 \\
3 & 0.0349 & 0.0361 & 0.0384 & 0.0401 & 0.0419 & 0.0436 & 0.0454 & 0.0646 & 0.0489 & 0.0506 \\
4 & 0.0698 & 0.0341 & 0.0558 & 0.0576 & 0.0593 & 0.0611 & 0.0628 & 0.0646 & 0.0663 & 0.0681 \\
\hline
5 & 0.0698 & 0.0716 & 0.0733 & 0.0750 & 0.0768 & 0.0785 & 0.0803 & 0.0820 & 0.1012 & 0.0855 \\
6 & 0.0873 & 0.0890 & 0.0908 & 0.0925 & 0.0942 & 0.0960 & 0.0977 & 0.0995 & 0.1012 & 0.1030 \\
7 & 0.1047 & 0.1065 & 0.1082 & 0.1100 & 0.1117 & 0.1134 & 0.1152 & 0.1169 & 0.1187 & 0.1204 \\
8 & 0.1222 & 0.1239 & 0.1257 & 0.1274 & 0.1292 & 0.1309 & 0.1326 & 0.1344 & 0.1361 & 0.1379 \\
9 & 0.1396 & 0.1414 & 0.1431 & 0.1449 & 0.1466 & 0.1484 & 0.1501 & 0.1518 & 0.1536 & 0.1553 \\
\hline
10 & 0.1571 & 0.1588 & 0.1606 & 0.1623 & 0.1815 & 0.1658 & 0.1675 & 0.1693 & 0.1710 & 0.1902 \\
11 & 0.1745 & 0.1763 & 0.1780 & 0.1798 & 0.1815 & 0.1833 & 0.1850 & 0.1867 & 0.1885 & 0.1902 \\
12 & 0.1920 & 0.1937 & 0.1955 & 0.1972 & 0.2164 & 0.2007 & 0.2025 & 0.2042 & 0.2059 & 0.2077 \\
13 & 0.2094 & 0.2286 & 0.2129 & 0.2147 & 0.2164 & 0.2182 & 0.2199 & 0.2217 & 0.2234 & 0.2251 \\
14 & 0.2269 & 0.2286 & 0.2304 & 0.2321 & 0.2339 & 0.2356 & 0.2374 & 0.2391 & 0.2409 & 0.2426 \\
\hline
15 & 0.2443 & 0.2461 & 0.2478 & 0.2496 & 0.2513 & 0.2531 & 0.2548 & 0.2566 & 0.2583 & 0.2600 \\
16 & 0.2618 & 0.2635 & 0.2827 & 0.2676 & 0.2688 & 0.2705 & 0.2723 & 0.2740 & 0.2758 & 0.2775 \\
17 & 0.2792 & 0.2810 & 0.2827 & 0.2845 & 0.2862 & 0.2880 & 0.2897 & 0.2915 & 0.2932 & 0.2950 \\
18 & 0.2967 & 0.2984 & 0.3002 & 0.3019 & 0.3211 & 0.3054 & 0.3072 & 0.3089 & 0.3107 & 0.3299 \\
19 & 0.3142 & 0.3159 & 0.3176 & 0.3194 & 0.3211 & 0.3229 & 0.3246 & 0.3264 & 0.3281 & 0.3299 \\
\hline
20 & 0.3316 & 0.3334 & 0.3351 & 0.3368 & 0.3386 & 0.3403 & 0.3421 & 0.3438 & 0.3456 & 0.3473 \\
21 & 0.3491 & 0.3508 & 0.3526 & 0.3543 & 0.3560 & 0.3578 & 0.3595 & 0.3613 & 0.3630 & 0.3648 \\
22 & 0.3665 & 0.3683 & 0.3875 & 0.3717 & 0.3735 & 0.3927 & 0.3944 & 0.3962 & 0.3805 & 0.3997 \\
23 & 0.3840 & 0.3857 & 0.3875 & 0.3892 & 0.3909 & 0.4102 & 0.3944 & 0.4136 & 0.3979 & 0.3997 \\
24 & 0.4014 & 0.4032 & 0.4049 & 0.4067 & 0.4084 & 0.4101 & 0.4294 & 0.4311 & 0.4154 & 0.4346 \\
\hline
25 & 0.4189 & 0.4206 & 0.4224 & 0.4241 & 0.4259 & 0.4451 & 0.4293 & 0.4485 & 0.4328 & 0.4346 \\
26 & 0.4363 & 0.4555 & 0.4398 & 0.4590 & 0.4433 & 0.4451 & 0.4468 & 0.4485 & 0.4503 & 0.4695 \\
27 & 0.4538 & 0.4555 & 0.4573 & 0.4590 & 0.4608 & 0.4625 & 0.4642 & 0.4660 & 0.4677 & 0.4695 \\
28 & 0.4712 & 0.4730 & 0.4922 & 0.4765 & 0.4782 & 0.4800 & 0.4817 & 0.4834 & 0.4852 & 0.4869 \\
29 & 0.4887 & 0.5079 & 0.4922 & 0.4939 & 0.4957 & 0.4974 & 0.4992 & 0.5009 & 0.5026 & 0.5044 \\
\hline
30 & 0.5061 & 0.5079 & 0.5096 & 0.5114 & 0.5131 & 0.5149 & 0.5166 & 0.5358 & 0.5376 & 0.5218 \\
31 & 0.5236 & 0.5253 & 0.5271 & 0.5288 & 0.5480 & 0.5323 & 0.5341 & 0.5358 & 0.5376 & 0.5393 \\
32 & 0.5410 & 0.5603 & 0.5445 & 0.5463 & 0.5480 & 0.5498 & 0.5515 & 0.5533 & 0.5550 & 0.5568 \\
33 & 0.5585 & 0.5602 & 0.5620 & 0.5637 & 0.5655 & 0.5672 & 0.5864 & 0.5882 & 0.5725 & 0.5917 \\
34 & 0.5759 & 0.5727 & 0.5794 & 0.5812 & 0.5829 & 0.5847 & 0.5864 & 0.5882 & 0.5899 & 0.6091 \\
\hline
35 & 0.5934 & 0.5951 & 0.5969 & 0.5986 & 0.6004 & 0.6021 & 0.6039 & 0.6056 & 0.6074 & 0.6091 \\
36 & 0.6109 & 0.6126 & 0.6143 & 0.6161 & 0.6178 & 0.6196 & 0.6213 & 0.6231 & 0.6248 & 0.6440 \\
37 & 0.6283 & 0.6301 & 0.6493 & 0.6335 & 0.6353 & 0.6370 & 0.6388 & 0.6405 & 0.6423 & 0.6440 \\
38 & 0.6458 & 0.6475 & 0.6667 & 0.6510 & 0.6527 & 0.6545 & 0.6562 & 0.6580 & 0.6597 & 0.6615 \\
39 & 0.6632 & 0.6650 & 0.6667 & 0.6684 & 0.6702 & 0.6719 & 0.6737 & 0.6754 & 0.6772 & 0.6964 \\
\hline
40 & 0.6807 & 0.6824 & 0.6842 & 0.6859 & 0.6876 & 0.7069 & 0.6911 & 0.7103 & 0.6946 & 0.6964 \\
41 & 0.6981 & 0.7173 & 0.7016 & 0.7034 & 0.7051 & 0.7243 & 0.7086 & 0.7278 & 0.7121 & 0.7138 \\
42 & 0.7156 & 0.7173 & 0.7365 & 0.7208 & 0.7400 & 0.7243 & 0.7260 & 0.7278 & 0.7470 & 0.7313 \\
43 & 0.7330 & 0.7348 & 0.7540 & 0.7383 & 0.7400 & 0.7418 & 0.7610 & 0.7452 & 0.7470 & 0.7487 \\
44 & 0.7505 & 0.7522 & 0.7540 & 0.7557 & 0.7575 & 0.7592 & 0.7784 & 0.7627 & 0.7644 & 0.7662 \\
\hline
45 & 0.7679 & 0.7871 & 0.7714 & 0.7906 & 0.7749 & 0.7767 & 0.7784 & 0.7801 & 0.7819 & 0.7836 \\
\hline
$\theta\text{ (degrees)}$ & 0 & 0.1 & 0.2 & 0.3 & 0.4 & 0.5 & 0.6 & 0.7 & 0.8 & 0.9 \\ \hline
\end{tabular}\par}
\clearpage

\subsection*{表 3.\quad （續）$360^\circ$ 與弧度換算表}
{\centering\fontsize{7.5pt}{8.8pt}\selectfont\setlength{\tabcolsep}{3.5pt}\renewcommand{\arraystretch}{0.96}
\begin{tabular}{|c|cccccccccc|}
\hline
$\theta\text{ (degrees)}$ & 0 & 0.1 & 0.2 & 0.3 & 0.4 & 0.5 & 0.6 & 0.7 & 0.8 & 0.9 \\ \hline
45 & 0.7854 & 0.7871 & 0.7889 & 0.7906 & 0.7924 & 0.7941 & 0.7959 & 0.7976 & 0.7994 & 0.8011 \\
46 & 0.7871 & 0.7889 & 0.7906 & 0.7924 & 0.7941 & 0.7959 & 0.7976 & 0.7993 & 0.8011 & 0.8186 \\
47 & 0.8046 & 0.8063 & 0.8081 & 0.8098 & 0.8116 & 0.8133 & 0.8151 & 0.8168 & 0.8185 & 0.8360 \\
48 & 0.8220 & 0.8238 & 0.8255 & 0.8273 & 0.8290 & 0.8308 & 0.8325 & 0.8343 & 0.8360 & 0.8535 \\
49 & 0.8395 & 0.8412 & 0.8430 & 0.8447 & 0.8465 & 0.8482 & 0.8500 & 0.8517 & 0.8535 & 0.8709 \\
\hline
50 & 0.8569 & 0.8587 & 0.8604 & 0.8622 & 0.8639 & 0.8657 & 0.8674 & 0.8692 & 0.8709 & 0.8884 \\
51 & 0.8744 & 0.8761 & 0.8779 & 0.8796 & 0.8814 & 0.8831 & 0.8849 & 0.8866 & 0.8884 & 0.9058 \\
52 & 0.8918 & 0.8936 & 0.8953 & 0.8971 & 0.8988 & 0.9006 & 0.9023 & 0.9041 & 0.9058 & 0.9233 \\
53 & 0.9093 & 0.9110 & 0.9128 & 0.9145 & 0.9163 & 0.9180 & 0.9198 & 0.9215 & 0.9233 & 0.9407 \\
54 & 0.9268 & 0.9285 & 0.9302 & 0.9320 & 0.9337 & 0.9355 & 0.9372 & 0.9390 & 0.9407 & 0.9582 \\
\hline
55 & 0.9442 & 0.9460 & 0.9477 & 0.9494 & 0.9512 & 0.9529 & 0.9547 & 0.9564 & 0.9582 & 0.9756 \\
56 & 0.9617 & 0.9634 & 0.9652 & 0.9669 & 0.9686 & 0.9704 & 0.9721 & 0.9739 & 0.9756 & 0.9931 \\
57 & 0.9791 & 0.9809 & 0.9983 & 0.9843 & 0.9861 & 0.9878 & 0.9896 & 0.9913 & 0.9931 & 1.0105 \\
58 & 0.9966 & 0.9983 & 1.0001 & 1.0018 & 1.0035 & 1.0053 & 1.0070 & 1.0088 & 1.0105 & 1.0280 \\
59 & 1.0140 & 1.0158 & 1.0175 & 1.0193 & 1.0210 & 1.0227 & 1.0245 & 1.0262 & 1.0280 & 1.0455 \\
\hline
60 & 1.0315 & 1.0332 & 1.0350 & 1.0367 & 1.0385 & 1.0402 & 1.0577 & 1.0437 & 1.0454 & 1.0629 \\
61 & 1.0489 & 1.0507 & 1.0524 & 1.0542 & 1.0559 & 1.0577 & 1.0594 & 1.0611 & 1.0629 & 1.0804 \\
62 & 1.0664 & 1.0681 & 1.0699 & 1.0716 & 1.0734 & 1.0751 & 1.0769 & 1.0786 & 1.0803 & 1.0978 \\
63 & 1.0838 & 1.0856 & 1.0873 & 1.0891 & 1.0908 & 1.0926 & 1.0943 & 1.0960 & 1.0978 & 1.1153 \\
64 & 1.1013 & 1.1030 & 1.1048 & 1.1065 & 1.1083 & 1.1100 & 1.1118 & 1.1135 & 1.1152 & 1.1327 \\
\hline
65 & 1.1345 & 1.1205 & 1.1222 & 1.1240 & 1.1257 & 1.1275 & 1.1292 & 1.1310 & 1.1327 & 1.1502 \\
66 & 1.1362 & 1.1379 & 1.1397 & 1.1414 & 1.1432 & 1.1449 & 1.1467 & 1.1484 & 1.1502 & 1.1676 \\
67 & 1.1536 & 1.1554 & 1.1571 & 1.1589 & 1.1606 & 1.1624 & 1.1641 & 1.1659 & 1.1676 & 1.1851 \\
68 & 1.1711 & 1.1728 & 1.1746 & 1.1763 & 1.1781 & 1.1798 & 1.1816 & 1.1833 & 1.1851 & 1.2025 \\
69 & 1.1885 & 1.1903 & 1.1920 & 1.1938 & 1.1955 & 1.1973 & 1.1990 & 1.2008 & 1.2025 & 1.2200 \\
\hline
70 & 1.2060 & 1.2077 & 1.2095 & 1.2112 & 1.2130 & 1.2147 & 1.2165 & 1.2182 & 1.2200 & 1.2374 \\
71 & 1.2235 & 1.2252 & 1.2269 & 1.2287 & 1.2304 & 1.2322 & 1.2339 & 1.2357 & 1.2374 & 1.2549 \\
72 & 1.2409 & 1.2427 & 1.2444 & 1.2461 & 1.2479 & 1.2496 & 1.2514 & 1.2531 & 1.2549 & 1.2723 \\
73 & 1.2584 & 1.2601 & 1.2619 & 1.2636 & 1.2811 & 1.2671 & 1.2688 & 1.2706 & 1.2723 & 1.2898 \\
74 & 1.2758 & 1.2776 & 1.2793 & 1.2811 & 1.2828 & 1.2845 & 1.2863 & 1.2880 & 1.2898 & 1.3073 \\
\hline
75 & 1.2933 & 1.3107 & 1.2968 & 1.2985 & 1.3002 & 1.3020 & 1.3037 & 1.3055 & 1.3072 & 1.3247 \\
76 & 1.3107 & 1.3125 & 1.3142 & 1.3160 & 1.3177 & 1.3194 & 1.3212 & 1.3229 & 1.3247 & 1.3422 \\
77 & 1.3282 & 1.3299 & 1.3317 & 1.3334 & 1.3352 & 1.3369 & 1.3386 & 1.3404 & 1.3421 & 1.3596 \\
78 & 1.3456 & 1.3474 & 1.3491 & 1.3509 & 1.3526 & 1.3701 & 1.3561 & 1.3578 & 1.3596 & 1.3771 \\
79 & 1.3631 & 1.3648 & 1.3666 & 1.3683 & 1.3701 & 1.3718 & 1.3736 & 1.3753 & 1.3770 & 1.3945 \\
\hline
80 & 1.3805 & 1.3823 & 1.3840 & 1.3858 & 1.3875 & 1.3893 & 1.3910 & 1.3927 & 1.3945 & 1.4120 \\
81 & 1.3980 & 1.3997 & 1.4015 & 1.4032 & 1.4050 & 1.4067 & 1.4085 & 1.4102 & 1.4119 & 1.4294 \\
82 & 1.4154 & 1.4172 & 1.4189 & 1.4207 & 1.4224 & 1.4242 & 1.4259 & 1.4277 & 1.4294 & 1.4469 \\
83 & 1.4329 & 1.4346 & 1.4364 & 1.4381 & 1.4399 & 1.4416 & 1.4434 & 1.4451 & 1.4469 & 1.4643 \\
84 & 1.4503 & 1.4521 & 1.4538 & 1.4556 & 1.4573 & 1.4591 & 1.4603 & 1.4626 & 1.4643 & 1.4818 \\
\hline
85 & 1.4678 & 1.4695 & 1.4713 & 1.4730 & 1.4748 & 1.4765 & 1.4783 & 1.4800 & 1.4818 & 1.4992 \\
86 & 1.4853 & 1.4870 & 1.4887 & 1.4905 & 1.4922 & 1.4940 & 1.4957 & 1.4975 & 1.4992 & 1.5167 \\
87 & 1.5027 & 1.5044 & 1.5062 & 1.5079 & 1.5097 & 1.5114 & 1.5132 & 1.5149 & 1.5167 & 1.5341 \\
88 & 1.5202 & 1.5219 & 1.5236 & 1.5254 & 1.5271 & 1.5289 & 1.5464 & 1.5324 & 1.5341 & 1.5516 \\
89 & 1.5376 & 1.5394 & 1.5411 & 1.5428 & 1.5446 & 1.5463 & 1.5481 & 1.5498 & 1.5516 & 1.5691 \\
\hline
90 & 1.5551 & 1.5568 & 1.5586 & 1.5603 & 1.5620 & 1.5638 & 1.5655 & 1.5673 & 1.5690 & 1.5865 \\
\hline
$\theta\text{ (degrees)}$ & 0 & 0.1 & 0.2 & 0.3 & 0.4 & 0.5 & 0.6 & 0.7 & 0.8 & 0.9 \\ \hline
\end{tabular}\par}


\end{document}
